\documentclass[11pt,a4paper]{article}

\usepackage[margin=0.85in]{geometry}
\usepackage{graphicx}
\usepackage{booktabs}
\usepackage{array}
\usepackage{float}
\usepackage{caption}
\usepackage{subcaption}
\usepackage{xcolor}
\usepackage{hyperref}
\usepackage{enumitem}
\usepackage{longtable}
\usepackage{fancyhdr}

\definecolor{accent}{HTML}{1F5A7A}
\definecolor{muted}{HTML}{5C6670}
\definecolor{lightgray}{HTML}{F3F5F7}

\hypersetup{
  colorlinks=true,
  linkcolor=accent,
  urlcolor=accent,
  citecolor=accent
}

\pagestyle{fancy}
\setlength{\headheight}{14pt}
\fancyhf{}
\lhead{\textcolor{muted}{SOM World Map}}
\rhead{\textcolor{muted}{Project Report}}
\cfoot{\thepage}

\setlength{\parskip}{0.55em}
\setlength{\parindent}{0pt}
\setlist[itemize]{leftmargin=1.35em, topsep=0.2em, itemsep=0.15em}
\setlist[enumerate]{leftmargin=1.35em, topsep=0.2em, itemsep=0.15em}

\newcommand{\figpath}{../../figures}
\newcommand{\datapath}{../../data}
\newcommand{\codepath}{../../src}

\title{
  \vspace{-1.0cm}
  \textbf{\Huge SOM World Map}\\
  \vspace{0.25cm}
  \Large Country Similarity, Temporal Trajectories, and SOM-Based Classification
}
\author{Neural Networks Final Project}
\date{May 2026}

\begin{document}

\maketitle

\begin{abstract}
This report summarizes a Self-Organizing Map (SOM) project built on World Bank country indicators. The project began as a static country similarity map and was extended into a temporal country-year SOM and a SOM-based majority-vote classifier. The central idea is to convert multidimensional national development profiles into an interpretable two-dimensional similarity landscape. The final workflow includes data collection, feature engineering, static SOM visualization, dynamic 2014-2023 trajectory analysis, Thailand-focused neighbor drift, time-shifted similarity, and a held-out income-level classification experiment that uses SOM nodes directly as class decision regions.
\end{abstract}

\setcounter{tocdepth}{1}
\tableofcontents
\newpage

\section{Executive Summary}

The project answers a simple question: if each country is represented by multiple development indicators, which countries are similar, and how does that similarity change over time? A SOM is a useful model for this question because it preserves neighborhood structure while reducing high-dimensional feature vectors to a two-dimensional grid.

The final project contains three analytical layers:

\begin{enumerate}
  \item \textbf{Static SOM:} a one-time country similarity map using the latest available World Bank indicators plus 2018 CO2 per capita.
  \item \textbf{Dynamic SOM:} a 2014-2023 country-year SOM that shows how countries move across a development-profile landscape over time.
  \item \textbf{SOM classifier:} a 4-node SOM majority-vote classifier for World Bank income level, evaluated on 15 held-out countries.
\end{enumerate}

\begin{table}[H]
\centering
\caption{Main project results.}
\begin{tabular}{ll}
\toprule
\textbf{Component} & \textbf{Result} \\
\midrule
Static country metadata & 217 countries \\
Static SOM complete cases & 171 countries \\
Static SOM grid & 9 x 10 \\
Static quantization error & 0.687 \\
Static topographic error & 0.018 \\
Dynamic country-year rows & 2,170 rows \\
Dynamic complete country-year rows & 1,733 rows \\
Countries complete for all 10 dynamic years & 158 countries \\
Dynamic SOM grid & 10 x 12 \\
Dynamic quantization error & 0.620 \\
Dynamic topographic error & 0.033 \\
SOM classifier split & 154 train countries, 15 test countries \\
SOM classifier accuracy & 0.800 \\
SOM classifier macro F1 & 0.774 \\
\bottomrule
\end{tabular}
\end{table}

\section{Data Sources and Feature Design}

The data source is the World Bank Indicators API. The static pipeline uses the latest available values from the 2018-2024 API window, then computes 2018 CO2 tons per capita from total CO2 emissions and population.

\subsection{Static SOM Features}

The static SOM uses seven features:

\begin{itemize}
  \item GDP per capita
  \item Life expectancy
  \item Internet users as a percentage of population
  \item Urban population as a percentage of total population
  \item Unemployment as a percentage of labor force
  \item Access to electricity as a percentage of population
  \item CO2 tons per capita in 2018
\end{itemize}

GDP and CO2 are log-transformed before scaling because their raw distributions are highly skewed. All model features are standardized before SOM training.

\begin{figure}[H]
\centering
\includegraphics[width=0.93\textwidth]{\figpath/feature_distributions.png}
\caption{Feature distributions after the log transformations used in the static SOM pipeline.}
\end{figure}

\subsection{Dynamic SOM Features}

The dynamic SOM uses country-year rows from 2014 through 2023. It trains on six high-coverage features:

\begin{itemize}
  \item GDP per capita
  \item Life expectancy
  \item Internet users as a percentage of population
  \item Urban population as a percentage of total population
  \item Unemployment as a percentage of labor force
  \item Access to electricity as a percentage of population
\end{itemize}

CO2 is still downloaded into the full country-year table where available, but it is not used as a 10-year SOM feature because the current endpoint does not provide complete recent coverage after 2018. This keeps the temporal SOM interpretable and avoids aggressive imputation.

\section{Static Country Similarity SOM}

The static SOM maps 171 complete country profiles onto a 9 x 10 grid. Countries with similar standardized feature vectors land near each other, even if they are geographically far apart.

\begin{figure}[H]
\centering
\includegraphics[width=0.86\textwidth]{\figpath/som_u_matrix_country_map.png}
\caption{Static SOM U-Matrix. Darker and lighter regions show changes in average neighbor distance across the SOM grid, while countries are placed at their best-matching units.}
\end{figure}

The U-Matrix gives a similarity landscape: dense regions indicate countries with comparable profiles, while high-distance boundaries suggest sharper transitions between profile groups.

\subsection{Feature Interpretation with Component Planes}

Component planes show how each standardized feature is distributed across the SOM grid. They make the model interpretable because each location on the grid can be read as a profile archetype.

\begin{figure}[H]
\centering
\includegraphics[width=0.93\textwidth]{\figpath/som_component_planes.png}
\caption{Static SOM component planes. Each panel shows the standardized SOM weight for one transformed feature.}
\end{figure}

\subsection{Macro Clusters}

To make the SOM easier to explain, K-Means is applied to SOM node weights and used to create five macro clusters. This is not a replacement for the SOM; it is a reporting layer on top of the continuous SOM landscape.

\begin{figure}[H]
\centering
\includegraphics[width=0.86\textwidth]{\figpath/som_macro_clusters.png}
\caption{Macro clusters overlaid on the static SOM grid.}
\end{figure}

\subsection{Thailand Similarity Case Study}

Thailand is used as a case study for nearest-neighbor interpretation. The static feature-space nearest countries to Thailand include Ecuador, Jamaica, Azerbaijan, Mexico, China, Dominican Republic, Belarus, Hungary, Uzbekistan, and Romania.

\begin{figure}[H]
\centering
\includegraphics[width=0.90\textwidth]{\figpath/similar_countries_THA.png}
\caption{Countries most similar to Thailand in scaled static feature space.}
\end{figure}

\section{Dynamic Country-Year SOM}

The dynamic extension changes the unit of analysis from a country to a country-year. A country such as Thailand becomes a sequence of points: Thailand 2014, Thailand 2015, and so on through Thailand 2023.

This allows the project to answer a stronger question: not only which countries are similar at one point in time, but also how national development profiles move across the similarity landscape.

\begin{table}[H]
\centering
\caption{Dynamic SOM data coverage.}
\begin{tabular}{lr}
\toprule
\textbf{Metric} & \textbf{Value} \\
\midrule
Year range & 2014-2023 \\
Total countries & 217 \\
Total country-year rows & 2,170 \\
Complete country-year rows for six core features & 1,733 \\
Countries with complete data for all 10 years & 158 \\
Countries represented in dynamic SOM model rows & 184 \\
\bottomrule
\end{tabular}
\end{table}

\subsection{Dynamic SOM Snapshots}

The dynamic SOM is trained on all complete country-year rows together. The following snapshots show how the distribution of country-year profiles changes over selected years.

\begin{figure}[H]
\centering
\includegraphics[width=0.96\textwidth]{\figpath/dynamic_som_year_snapshots.png}
\caption{Dynamic SOM snapshots for selected years. Each point is a country-year observation.}
\end{figure}

\begin{figure}[H]
\centering
\includegraphics[width=0.94\textwidth]{\figpath/dynamic_som_component_planes.png}
\caption{Dynamic SOM component planes for the six time-series features.}
\end{figure}

\subsection{Geographic View of Dynamic Clusters}

The dynamic SOM clusters can also be displayed on a geographic scatter plot. This view is not a replacement for a formal map, but it gives a quick visual check of how SOM profile clusters are distributed across the world.

\begin{figure}[H]
\centering
\includegraphics[width=0.96\textwidth]{\figpath/dynamic_world_cluster_scatter.png}
\caption{Geographic scatter of dynamic SOM macro clusters.}
\end{figure}

\subsection{Thailand Trajectory}

Thailand has complete dynamic data for all years from 2014 through 2023. Its trajectory shows how the country moves through the SOM similarity landscape as the underlying feature profile changes.

\begin{figure}[H]
\centering
\includegraphics[width=0.72\textwidth]{\figpath/dynamic_som_trajectory_THA.png}
\caption{Thailand trajectory across the dynamic SOM grid from 2014 to 2023.}
\end{figure}

The trajectory shows a movement from the 2014-2017 region into a later position after 2018. The feature changes behind this movement include rising internet adoption and urbanization, while GDP and unemployment fluctuate across the pandemic period.

\subsection{Neighbor Drift}

Neighbor drift records the closest country to Thailand for each year in standardized feature space. The closest peer changes over time:

\begin{table}[H]
\centering
\caption{Thailand's nearest dynamic peer by year.}
\begin{tabular}{rlll}
\toprule
\textbf{Year} & \textbf{Nearest peer} & \textbf{Country code} & \textbf{Distance} \\
\midrule
2014 & Viet Nam & VNM & 0.946 \\
2015 & Ecuador & ECU & 0.944 \\
2016 & China & CHN & 0.876 \\
2017 & Ecuador & ECU & 0.812 \\
2018 & Ecuador & ECU & 0.608 \\
2019 & Ecuador & ECU & 0.680 \\
2020 & Poland & POL & 0.724 \\
2021 & Poland & POL & 0.814 \\
2022 & Jamaica & JAM & 0.722 \\
2023 & Jamaica & JAM & 0.743 \\
\bottomrule
\end{tabular}
\end{table}

\begin{figure}[H]
\centering
\includegraphics[width=0.86\textwidth]{\figpath/neighbor_drift_THA.png}
\caption{Distance from Thailand to its closest country-year peer by year. Labels show the nearest peer country code.}
\end{figure}

\subsection{Movement Scores}

Movement score is calculated as cumulative step length in standardized feature space over the observed time window. For reporting, complete-window countries are prioritized so that partial-window outliers do not dominate the narrative.

\begin{figure}[H]
\centering
\includegraphics[width=0.86\textwidth]{\figpath/country_movement_scores.png}
\caption{Countries with the largest complete-window 2014-2023 profile movement.}
\end{figure}

\subsection{Time-Shifted Similarity}

Time-shifted similarity compares Thailand in 2023 to every other country-year in the dynamic dataset. This supports a stronger interpretation than ordinary nearest-neighbor search because it can identify historical country profiles similar to Thailand's current profile.

\begin{figure}[H]
\centering
\includegraphics[width=0.86\textwidth]{\figpath/time_shifted_similarity_THA_2023.png}
\caption{Most similar country-years to Thailand in 2023.}
\end{figure}

\section{SOM-Based Classification Add-On}

The project also includes a supervised validation layer, but the classifier is still SOM-based. The goal is not to train a separate black-box classifier. Instead, the SOM itself is used as the classification surface.

\subsection{Classification Task}

The target label is World Bank income level:

\begin{itemize}
  \item Low income
  \item Lower middle income
  \item Upper middle income
  \item High income
\end{itemize}

Countries labeled ``Not classified'' are removed. The classifier uses the same static feature set as the static SOM. The train-test split contains 154 training countries and 15 held-out test countries.

\subsection{Method}

The classifier uses a 4-node SOM, matching the four income classes. After training on the train set, each SOM node is assigned the majority income class among the training countries mapped to that node. A held-out test country is classified by finding its best-matching SOM node and using that node's majority label.

This design keeps the method aligned with the project theme: classification is interpreted through the SOM grid itself.

\begin{figure}[H]
\centering
\includegraphics[width=0.74\textwidth]{\figpath/som_income_classifier_map.png}
\caption{SOM majority-vote income classifier. Background node colors show majority class labels from the train set. Test countries are overlaid on their best-matching nodes.}
\end{figure}

\subsection{Classifier Results}

The SOM classifier achieved 0.800 accuracy and 0.774 macro F1 on 15 held-out countries.

\begin{table}[H]
\centering
\caption{SOM income classifier test predictions.}
\begin{tabular}{llll}
\toprule
\textbf{Country} & \textbf{Actual} & \textbf{Predicted} & \textbf{Correct} \\
\midrule
GNB & Low & Low & Yes \\
PER & Upper middle & Upper middle & Yes \\
PNG & Lower middle & Low & No \\
MWI & Low & Low & Yes \\
SLB & Lower middle & Low & No \\
ISR & High & High & Yes \\
DJI & Lower middle & Lower middle & Yes \\
SWE & High & High & Yes \\
MUS & Upper middle & Lower middle & No \\
IRQ & Upper middle & Upper middle & Yes \\
POL & High & High & Yes \\
NOR & High & High & Yes \\
CZE & High & High & Yes \\
IRN & Upper middle & Upper middle & Yes \\
VNM & Lower middle & Lower middle & Yes \\
\bottomrule
\end{tabular}
\end{table}

\begin{figure}[H]
\centering
\includegraphics[width=0.72\textwidth]{\figpath/som_income_classifier_confusion_matrix.png}
\caption{Confusion matrix for the SOM income classifier.}
\end{figure}

\section{Repository Outputs}

The main reproducible scripts are:

\begin{itemize}
  \item \texttt{src/som\_world\_map\_pipeline.py}: static SOM data, model outputs, and figures.
  \item \texttt{src/dynamic\_som\_pipeline.py}: 2014-2023 dynamic country-year SOM.
  \item \texttt{src/som\_income\_classifier.py}: SOM majority-vote income classifier.
\end{itemize}

\begin{longtable}{p{0.36\textwidth}p{0.54\textwidth}}
\caption{Key generated artifacts.}\\
\toprule
\textbf{Artifact} & \textbf{Purpose} \\
\midrule
\endfirsthead
\toprule
\textbf{Artifact} & \textbf{Purpose} \\
\midrule
\endhead
\path{data/country_som_results.csv} & Static SOM country positions and macro clusters. \\
\path{data/thailand_similar_countries.csv} & Static nearest-neighbor table for Thailand. \\
\path{data/country_year_som_results.csv} & Dynamic country-year SOM positions. \\
\path{data/neighbor_drift_THA.csv} & Thailand nearest peers by year. \\
\path{data/movement_scores.csv} & Country profile movement scores over time. \\
\path{data/time_shifted_similarity_THA_2023.csv} & Country-years most similar to Thailand in 2023. \\
\path{data/som_income_classifier_metrics.json} & SOM classifier metrics and test predictions. \\
\path{figures/som_u_matrix_country_map.png} & Static SOM U-Matrix visualization. \\
\path{figures/dynamic_som_trajectory_THA.png} & Thailand dynamic SOM trajectory. \\
\path{figures/som_income_classifier_map.png} & SOM majority-vote classification map. \\
\bottomrule
\end{longtable}

\section{Limitations}

\begin{itemize}
  \item The SOM is an exploratory model. It reveals similarity structure but does not prove causality.
  \item Feature selection strongly affects the map. Countries are similar only with respect to the selected indicators.
  \item Static CO2 per capita is computed for 2018 because the current World Bank endpoint does not provide complete recent coverage.
  \item The dynamic SOM uses complete cases for six core features. This avoids hidden imputation assumptions but reduces the number of usable country-year rows.
  \item World Bank income level is a useful classification target, but the classifier should be interpreted as a development-profile validation experiment, not a production income classifier.
  \item The SOM classifier uses only 15 held-out test countries, so the accuracy should be presented as a small validation result rather than a definitive benchmark.
\end{itemize}

\section{Recommended Next Steps}

The project is now strong enough to present as a prototype with three connected layers: static similarity, temporal movement, and SOM-based classification. The next improvements should focus on polish rather than adding unnecessary model complexity.

\begin{enumerate}
  \item Build a simple Streamlit dashboard with tabs for static SOM, dynamic trajectory, neighbor drift, and classifier results.
  \item Add a country selector so the trajectory and neighbor drift analysis can be repeated for any country, not only Thailand.
  \item Add a formal train-test robustness check for the SOM classifier with repeated stratified splits.
  \item Add a short appendix explaining why CO2 is excluded from the dynamic SOM feature set.
  \item Replace the geographic scatter with a true choropleth map if a dashboard is built.
\end{enumerate}

\section{Conclusion}

The final project is no longer just a basic SOM demonstration. It is a compact analytical system for visualizing national development profiles. The static SOM provides an interpretable country similarity map, the dynamic SOM shows how countries move over a decade, and the SOM-based classifier demonstrates that the learned similarity landscape can also support a direct majority-vote classification task. The project remains technically simple enough to explain clearly, while adding enough temporal and supervised structure to feel substantially more complete than a basic prototype.

\end{document}
